\documentclass[10pt,a4paper]{amsart}
\usepackage{amsmath,amssymb,amsthm,graphicx}
\usepackage{verbatim}
\newenvironment{reference}{\comment}{\endcomment}
%\newenvironment{reference}{}{}
\newenvironment{slogan}{\comment}{\endcomment}
\newenvironment{history}{\comment}{\endcomment}

% For commutative diagrams we use Xy-pic
\usepackage[all]{xy}

% We use 2cell for 2-commutative diagrams.
\xyoption{2cell}
% Theorem environments.
%
\theoremstyle{plain}
\newtheorem{theorem}[subsection]{Theorem}
\newtheorem{proposition}[subsection]{Proposition}
\newtheorem{lemma}[subsection]{Lemma}

\theoremstyle{definition}
\newtheorem{definition}[subsection]{Definition}
\newtheorem{example}[subsection]{Example}
\newtheorem{exercise}[subsection]{Exercise}
\newtheorem{situation}[subsection]{Situation}

\theoremstyle{remark}
\newtheorem{remark}[subsection]{Remark}
\newtheorem{remarks}[subsection]{Remarks}

\numberwithin{equation}{subsection}

% Macros
%
\def\lim{\mathop{\mathrm{lim}}\nolimits}
\def\colim{\mathop{\mathrm{colim}}\nolimits}
\def\Spec{\mathop{\mathrm{Spec}}}
\def\Hom{\mathop{\mathrm{Hom}}\nolimits}
\def\Ext{\mathop{\mathrm{Ext}}\nolimits}
\def\SheafHom{\mathop{\mathcal{H}\!\mathit{om}}\nolimits}
\def\SheafExt{\mathop{\mathcal{E}\!\mathit{xt}}\nolimits}
\def\Sch{\mathit{Sch}}
\def\Mor{\mathop{\mathrm{Mor}}\nolimits}
\def\Ob{\mathop{\mathrm{Ob}}\nolimits}
\def\Sh{\mathop{\mathit{Sh}}\nolimits}
\def\NL{\mathop{N\!L}\nolimits}
\def\CH{\mathop{\mathrm{CH}}\nolimits}
\def\proetale{{pro\text{-}\acute{e}tale}}
\def\etale{{\acute{e}tale}}
\def\QCoh{\mathit{QCoh}}
\def\Ker{\mathop{\mathrm{Ker}}}
\def\Im{\mathop{\mathrm{Im}}}
\def\Coker{\mathop{\mathrm{Coker}}}
\def\Coim{\mathop{\mathrm{Coim}}}

% Boxtimes
%
\DeclareMathSymbol{\boxtimes}{\mathbin}{AMSa}{"02}

%
% Macros for moduli stacks/spaces
%
\def\QCohstack{\mathcal{QC}\!\mathit{oh}}
\def\Cohstack{\mathcal{C}\!\mathit{oh}}
\def\Spacesstack{\mathcal{S}\!\mathit{paces}}
\def\Quotfunctor{\mathrm{Quot}}
\def\Hilbfunctor{\mathrm{Hilb}}
\def\Curvesstack{\mathcal{C}\!\mathit{urves}}
\def\Polarizedstack{\mathcal{P}\!\mathit{olarized}}
\def\Complexesstack{\mathcal{C}\!\mathit{omplexes}}
% \Pic is the operator that assigns to X its picard group, usage \Pic(X)
% \Picardstack_{X/B} denotes the Picard stack of X over B
% \Picardfunctor_{X/B} denotes the Picard functor of X over B
\def\Pic{\mathop{\mathrm{Pic}}\nolimits}
\def\Picardstack{\mathcal{P}\!\mathit{ic}}
\def\Picardfunctor{\mathrm{Pic}}
\def\Deformationcategory{\mathcal{D}\!\mathit{ef}}
\usepackage{cleveref}
\crefname{originaltheorem}{original theorem}{original theorems}
\Crefname{originaltheorem}{Original theorem}{Original theorems}
\makeatletter
\newlabel{etale-theorem-etale-radicial-open}{{025G}{source}{}{sourcepointer-1}{}}
\newlabel{etale-theorem-etale-radicial-open@cref}{{[originaltheorem][0][]025G}{[1][0][]source}{}{}{}}
\makeatother
\crefname{originallemma}{original lemma}{original lemmas}
\Crefname{originallemma}{Original lemma}{Original lemmas}
\makeatletter
\newlabel{lemma-properties-diagonal}{{03HK}{source}{}{sourcepointer-2}{}}
\newlabel{lemma-properties-diagonal@cref}{{[originallemma][0][]03HK}{[1][0][]source}{}{}{}}
\makeatother
\crefname{originallemma}{original lemma}{original lemmas}
\Crefname{originallemma}{Original lemma}{Original lemmas}
\makeatletter
\newlabel{morphisms-lemma-diagonal-unramified-morphism}{{02GE}{source}{}{sourcepointer-3}{}}
\newlabel{morphisms-lemma-diagonal-unramified-morphism@cref}{{[originallemma][0][]02GE}{[1][0][]source}{}{}{}}
\makeatother
\crefname{originallemma}{original lemma}{original lemmas}
\Crefname{originallemma}{Original lemma}{Original lemmas}
\makeatletter
\newlabel{schemes-lemma-diagonal-immersion}{{01KJ}{source}{}{sourcepointer-4}{}}
\newlabel{schemes-lemma-diagonal-immersion@cref}{{[originallemma][0][]01KJ}{[1][0][]source}{}{}{}}
\makeatother
\crefname{originallemma}{original lemma}{original lemmas}
\Crefname{originallemma}{Original lemma}{Original lemmas}
\makeatletter
\newlabel{spaces-lemma-representable-transformations-property-implication}{{02YO}{source}{}{sourcepointer-5}{}}
\newlabel{spaces-lemma-representable-transformations-property-implication@cref}{{[originallemma][0][]02YO}{[1][0][]source}{}{}{}}
\makeatother
\crefname{originallemma}{original lemma}{original lemmas}
\Crefname{originallemma}{Original lemma}{Original lemmas}
\makeatletter
\newlabel{spaces-properties-lemma-etale-local}{{03FS}{source}{}{sourcepointer-6}{}}
\newlabel{spaces-properties-lemma-etale-local@cref}{{[originallemma][0][]03FS}{[1][0][]source}{}{}{}}
\makeatother
\crefname{originallemma}{original lemma}{original lemmas}
\Crefname{originallemma}{Original lemma}{Original lemmas}
\makeatletter
\newlabel{spaces-properties-lemma-etale-permanence}{{03FV}{source}{}{sourcepointer-7}{}}
\newlabel{spaces-properties-lemma-etale-permanence@cref}{{[originallemma][0][]03FV}{[1][0][]source}{}{}{}}
\makeatother
\title{A morphism of algebraic spaces is unramified exactly when its diagonal is an open immersion}
\author{The Stacks Project Authors}
\date{stacks-project:spaces-morphisms.tex:local-snapshot-20261012; GNU FDL 1.2 or later}
\begin{document}
\maketitle
\setcounter{subsection}{1}
\noindent\textit{Editorial note.} All delivered mathematical text is the author's own, copied from the pinned Stacks Project chapter file spaces-morphisms.tex; no statement or proof step was rewritten, reordered or removed.\par\medskip
\noindent\textit{Editorial note.} Both directions of the equivalence are written out by the author, with the cartesian diagram as the working object; the local finite type hypothesis in the second direction is part of the statement span.\par\medskip
\noindent\textit{Editorial note.} The seven references of the counted spans are recorded as non-load-bearing navigation pointers; the argument that consumes them stays inside the counted proof span.\par\medskip
\noindent\textit{Editorial note.} Preamble: the three native excerpt spans of sources/stacks-spaces-morphisms/preamble.tex (lines 8-13, 14-18 and 154-217). The chapter file itself carries no preamble text.\par\medskip
\noindent\textit{Allowed standard prerequisites (editorial).}\begin{itemize}
\item The diagonal of a morphism of algebraic spaces is a representable monomorphism. Lemma lemma-properties-diagonal, used at proof lines 7737-7738. Supplies the input that has to be upgraded to an open immersion.
\item Etaleness is local on the source and is preserved by etale base change of schemes over Y. Properties of Spaces, Lemmas spaces-properties-lemma-etale-local and spaces-properties-lemma-etale-permanence, used at proof lines 7752-7767 and 7780-7784. Transfers the property of the diagonal between the scheme and the algebraic space picture.
\item The scheme case: unramified equals open diagonal. Morphisms, Lemma morphisms-lemma-diagonal-unramified-morphism, used at proof lines 7761-7763 and 7789-7790. Is the input applied to U -> V in both directions.
\item A representable etale monomorphism is an open immersion. Etale Morphisms, Theorem etale-theorem-etale-radicial-open, with Spaces, Lemma spaces-lemma-representable-transformations-property-implication, used at proof lines 7768-7774. Converts the etale monomorphism into an open immersion.
\item The diagonal of a morphism of schemes is an immersion. Schemes, Lemma schemes-lemma-diagonal-immersion, used at proof lines 7785-7787. Provides the immersion half in the converse direction.
\end{itemize}\medskip
\ifdefined\hypertarget\hypertarget{sourcepointer-1}{}\fi
\noindent\textit{Source reference.} Original theorem 025G: focused-seeds/AGP/sources/stacks-snapshot-20261011/tags.txt, source lines 2813--2813. Navigation pointer into the Stacks Project; the standard fact it supplies is applied inside the counted proof span and the argument resting on it stays there.\par\medskip
\ifdefined\hypertarget\hypertarget{sourcepointer-2}{}\fi
\noindent\textit{Source reference.} Original lemma 03HK: focused-seeds/AGP/sources/stacks-snapshot-20261011/tags.txt, source lines 4570--4570. Navigation pointer into the Stacks Project; the standard fact it supplies is applied inside the counted proof span and the argument resting on it stays there.\par\medskip
\ifdefined\hypertarget\hypertarget{sourcepointer-3}{}\fi
\noindent\textit{Source reference.} Original lemma 02GE: focused-seeds/AGP/sources/stacks-snapshot-20261011/tags.txt, source lines 3225--3225. Navigation pointer into the Stacks Project; the standard fact it supplies is applied inside the counted proof span and the argument resting on it stays there.\par\medskip
\ifdefined\hypertarget\hypertarget{sourcepointer-4}{}\fi
\noindent\textit{Source reference.} Original lemma 01KJ: focused-seeds/AGP/sources/stacks-snapshot-20261011/tags.txt, source lines 2052--2052. Navigation pointer into the Stacks Project; the standard fact it supplies is applied inside the counted proof span and the argument resting on it stays there.\par\medskip
\ifdefined\hypertarget\hypertarget{sourcepointer-5}{}\fi
\noindent\textit{Source reference.} Original lemma 02YO: focused-seeds/AGP/sources/stacks-snapshot-20261011/tags.txt, source lines 3883--3883. Navigation pointer into the Stacks Project; the standard fact it supplies is applied inside the counted proof span and the argument resting on it stays there.\par\medskip
\ifdefined\hypertarget\hypertarget{sourcepointer-6}{}\fi
\noindent\textit{Source reference.} Original lemma 03FS: focused-seeds/AGP/sources/stacks-snapshot-20261011/tags.txt, source lines 4503--4503. Navigation pointer into the Stacks Project; the standard fact it supplies is applied inside the counted proof span and the argument resting on it stays there.\par\medskip
\ifdefined\hypertarget\hypertarget{sourcepointer-7}{}\fi
\noindent\textit{Source reference.} Original lemma 03FV: focused-seeds/AGP/sources/stacks-snapshot-20261011/tags.txt, source lines 4506--4506. Navigation pointer into the Stacks Project; the standard fact it supplies is applied inside the counted proof span and the argument resting on it stays there.\par\medskip
\begin{lemma}
\label{lemma-diagonal-unramified-morphism}
Let $S$ be a scheme.
Let $f : X \to Y$ be a morphism of algebraic spaces over $S$.
\begin{enumerate}
\item If $f$ is unramified, then the diagonal morphism
$\Delta_{X/Y} : X \to X \times_Y X$ is an open immersion.
\item If $f$ is locally of finite type
and $\Delta_{X/Y}$ is an open immersion, then $f$ is unramified.
\end{enumerate}
\end{lemma}

\begin{proof}
We know in any case that $\Delta_{X/Y}$ is a representable monomorphism, see
Lemma \ref{lemma-properties-diagonal}.
Choose a scheme $V$ and a surjective \'etale morphism $V \to Y$.
Choose a scheme $U$ and a surjective \'etale morphism $U \to X \times_Y V$.
Consider the commutative diagram
$$
\xymatrix{
U \ar[d] \ar[rr]_-{\Delta_{U/V}} & &
U \times_V U \ar[d] \ar[r] &
V \ar[d]^{\Delta_{V/Y}} \\
X \ar[rr]^-{\Delta_{X/Y}} & &
X \times_Y X \ar[r] &
V \times_Y V
}
$$
with cartesian right square. The left vertical arrow is surjective \'etale.
The right vertical arrow is \'etale as a morphism between schemes
\'etale over $Y$, see
Properties of Spaces,
Lemma \ref{spaces-properties-lemma-etale-permanence}.
Hence the middle vertical arrow is \'etale too (but it need not be
surjective).

\medskip\noindent
Assume $f$ is unramified. Then $U \to V$ is unramified, hence
$\Delta_{U/V}$ is an open immersion by
Morphisms, Lemma \ref{morphisms-lemma-diagonal-unramified-morphism}.
Looking at the left square of the diagram above we conclude that
$\Delta_{X/Y}$ is an \'etale morphism, see
Properties of Spaces,
Lemma \ref{spaces-properties-lemma-etale-local}.
Hence $\Delta_{X/Y}$ is a representable \'etale monomorphism, which
implies that it is an open immersion by
\'Etale Morphisms, Theorem \ref{etale-theorem-etale-radicial-open}.
(See also
Spaces, Lemma
\ref{spaces-lemma-representable-transformations-property-implication}
for the translation from schemes language into the language of functors.)

\medskip\noindent
Assume that $f$ is locally of finite type and that $\Delta_{X/Y}$
is an open immersion. This implies that $U \to V$ is locally of finite
type too (by definition of a morphism of algebraic spaces which is
locally of finite type). Looking at the displayed diagram above
we conclude that $\Delta_{U/V}$ is \'etale as a morphism between
schemes \'etale over $X \times_Y X$, see
Properties of Spaces,
Lemma \ref{spaces-properties-lemma-etale-permanence}.
But since $\Delta_{U/V}$ is the diagonal of a morphism between schemes
we see that it is in any case an immersion, see
Schemes, Lemma \ref{schemes-lemma-diagonal-immersion}.
Hence it is an open immersion, and we conclude
that $U \to V$ is unramified by
Morphisms, Lemma \ref{morphisms-lemma-diagonal-unramified-morphism}.
This in turn means that $f$ is unramified by definition.
\end{proof}

\end{document}
